
In this work, we present a limited exploration of the SCD pretraining objective across multiple domains of pretraining data and property prediction tasks.  We find that SCD pretraining provides robust improvements on all tasks investigated and matches or outperforms the benefits of force-energy pretraining, as well as other pretraining methods. Small SCD pretrained models achieve state-of-the-art results across multiple domains: QM9,  Matbench properties, and Ligand binding affinity benchmarks. Overall, our findings can be summarized through several key observations:

\textbf{All data is useful for pretraining:} We observed that SCD pretraining on atomistic data from any domain or source (DFT accurate ground-state, 'non-equilibrium', or diffusion-generated) conveyed an improvement across all benchmarking tasks. Surprisingly, we find that even ligand binding affinity predictions are improved when pretraining on periodic materials (+12\%), and pretraining with non-periodic, non-equilibrium molecule geometries (OMol25) conveyed a surprising benefit to bandgap predictions in periodic materials (+26.9\%) (see Table \ref{results_summary}). However, we hypothesize that at least some of these benefits may stem from pre-smoothing of the representation space prior to task-specific fine-tuning (Figure \ref{fig:umap}).

\textbf{Pretraining is not necessarily expensive:}
Our backbone architecture is neither complex nor state-of-the-art, and without pretraining, it provides only a modest baseline. However, we find that a simple, fast GNN pretrained by SCD can outperform supervised learning with advanced tensor-product based architectures (e.g., EquiformerV2 \cite{EquiformerV2_Liao2023}) in both prediction accuracy and total compute cost (including pretraining GPU-hours. Table \ref{tab:scd_vs_jmp_qm9}).  Furthermore, large datasets are not always necessary to realize these benefits. Pretraining on a random 10\% slice of the PCQ dataset (300k molecules) captures 97\% of the performance gain of the full dataset for QM9 tasks, and pretraining on just 607k materials yields band gap accuracies within 1.7\% of larger models pretrained on 120M force-energy labeled geometries, 108M of which are from materials. Overall, we find that efficient pretraining is more impactful and cost effective than implementing inductive bias through complex architectural elements (as has been found in other fields). 

\textbf{Data relevance is key to performance:}
Across all tasks explored, the strongest predictor of success was similarity between samples in the pretraining dataset and the fine-tuning dataset. This was most apparent in materials bandgap prediction, in which pretraining datasets containing inorganic atoms provided far greater benefit than those containing only organic elements. Notably, our model pretrained on a combine super-set of single domain data (CT-SCD-ALL) achieved significantly better average performance across all domains, but did not surpass the performance of copies pretrained on single-domain data for in-domain tasks. We attribute this small drop in performance to the limited expressivity of our 10M parameter model rather than representation conflicts in multi-domain training, though further studies are needed to confirm this. We expect larger models to receive greater benefit from pretraining with large multi-domain datasets.


\textbf{Pretraining a Foundation Model does not require labels:}
Force-energy pretraining has proven to be an effective and scalable pretraining strategy when DFT labels are readily available \cite{JMP_Shoghi2023, HackNIP_kim2025}. However, computing large numbers of new ground truth (DFT) forces and energies is extremely expensive and time-consuming. Finding effective methods of pretraining with unlabeled, low fidelity, and or cheaply generated atomistic geometries could enable further generalization of atomistic foundation models without requiring a significant expansion of DFT labeled datasets. We provide SCD as a demonstration that this is possible through self-supervised learning. When controlling for architecture and pretraining geometries (labeled vs unlabeled), SCD pretraining can match or surpass the performance of supervised force-energy pretraining in downstream property prediction tasks. Further, we show that even diffusion generated conformers (i.e. SAIR) benefit SCD pretraining. As generative methods in the atomistic sciences continue to advance, we anticipate an increased availability of large 'synthetic' datasets for pretraining. 

In developing SCD, our aim is to inspire further exploration into self-supervised learning and efficient pretraining for atomistic foundation models. More work is still needed to investigate the effects of scaling across dataset and model dimensions. To assist in this effort, we've made our code and pre-trained models publicly available at \url{https://github.com/TyJPerez/SelfConditionedDenoisingAtoms}

% intend to make our code, including dataset loaders, training algorithms, and pretrained model weights, available online.
% % TODO Add links to github, etc

% Our code and pretrained models are publicly available at \url{https://github.com/TyJPerez/SelfConditionedDenoisingAtoms}



